% R_V Paper - COLM 2026 Submission
% "Self-Referential Processing Induces Geometric Contraction in Transformer Value Spaces"
% Abstract deadline: March 26, 2026
% Full paper deadline: March 31, 2026

\documentclass{article}

% COLM 2026 style (replace with official when available)
\usepackage[utf8]{inputenc}
\usepackage[T1]{fontenc}
\usepackage{graphicx}
\usepackage{amsmath,amssymb,amsfonts}
\usepackage{booktabs}
\usepackage{natbib}
\usepackage{hyperref}
\usepackage{microtype}
\usepackage{xcolor}
\usepackage{multirow}
\usepackage{subcaption}

% Page setup
\usepackage[margin=1in]{geometry}

% Graphics path
\graphicspath{{figures/}}

% Custom commands
\newcommand{\RV}{R_V}
\newcommand{\PR}{\text{PR}}
\newcommand{\deff}{d_{\text{eff}}}

\title{Self-Referential Processing Induces Geometric Contraction\\in Transformer Value Spaces}

\author{
  Anonymous Authors\\
  % Affiliation suppressed for review
}

\date{}

\begin{document}

\maketitle

%% ============================================================
%% ABSTRACT (~250 words)
%% ============================================================

\begin{abstract}
We investigate the geometric signatures of self-referential processing in transformer value spaces using a novel metric, the relative participation ratio ($\RV$).
Across five architectures (Mistral-7B, OPT-6.7B, GPT-2 XL, Qwen2.5-7B, Pythia-1.4B) and 10 computational modes, we find that self-referential prompts induce consistent geometric contraction ($\RV < 1$) localized to late integration layers.
The effect is robust: 30/36 tests survive FDR correction ($\alpha=0.05$), 10/13 core effects survive conservative cluster-robust standard errors (DEFF=2), and the signature persists after controlling for perplexity confounds ($d=-1.80$, $p<10^{-10}$).

A full head sweep across all 1,024 attention heads in Mistral-7B reveals 606 heads with significant $\RV$ separation, with the strongest effects in middle layers (L10H20, $d=3.90$).
SVD circuit decomposition identifies rank-deficiency patterns specific to self-referential processing.
Linear probes achieve 100\% classification accuracy from Layer 4 onward, and representational similarity analysis reveals a distinctive distance trajectory for self-referential representations across layers.

For AI safety applications, $\RV$ achieves AUROC=0.909 for detecting self-referential content, but critically tracks \emph{content} rather than \emph{intent}: genuinely and deceptively self-referential text are geometrically indistinguishable ($d=-0.06$).
All analyses use frequentist and Bayesian inference with pre-registered effect sizes, power analysis, and multiple comparison corrections.
We release all code, prompts, and computed statistics.
\end{abstract}

%% ============================================================
%% 1. INTRODUCTION
%% ============================================================

\section{Introduction}

% Hook: why self-referential processing matters for interpretability
% Gap: no geometric characterization of how transformers process self-reference
% Contribution: R_V metric, cross-architecture validation, causal evidence, safety implications

% Key contributions (6):
% 1. R_V metric definition and spectral fingerprinting across 10 modes
% 2. Cross-architecture universality (5 models)
% 3. Causal validation: necessity (d=3.29) and KV sufficiency (OR=13.96)
% 4. Circuit-level decomposition: 1024-head sweep + SVD
% 5. Statistical hardening: FDR, cluster-robust SEs, perplexity controls
% 6. Safety: AUROC=0.909, content-not-intent finding

\textbf{TODO: Write introduction.}


%% ============================================================
%% 2. RELATED WORK
%% ============================================================

\section{Related Work}

% Mechanistic interpretability: circuits, probing, activation patching
% Geometric analysis of representations: manifold hypothesis, effective rank
% Self-referential processing in LLMs: behavioral studies
% AI safety: deceptive alignment, monitoring proposals

\textbf{TODO: Write related work.}


%% ============================================================
%% 3. METHODS
%% ============================================================

\section{Methods}

\subsection{The Relative Participation Ratio ($\RV$)}

% Definition of PR from SVD
\begin{equation}
    \PR(L) = \frac{\left(\sum_i \sigma_i^2(L)\right)^2}{\sum_i \sigma_i^4(L)}
\end{equation}

% Definition of R_V
\begin{equation}
    \RV = \frac{\PR(L_{\text{late}})}{\PR(L_{\text{early}})}
    \label{eq:rv}
\end{equation}

% Interpretation: R_V < 1 = contraction, R_V > 1 = expansion

\subsection{Prompt Design and Controls}

% 10 computational modes (Table 1)
% Template structure: L5_refined, L4_full (recursive), baseline_creative, baseline_math
% N=300 per condition for main experiments
% Perplexity matching protocol

\subsection{Statistical Framework}

% Pre-registered pipeline:
% 1. Cohen's d with 95% CI
% 2. Welch's t-test (two-sided)
% 3. Bayes factors (BF10)
% 4. FDR correction (Benjamini-Hochberg)
% 5. Cluster-robust SEs (ICC + DEFF)
% 6. Achieved power analysis

\subsection{Models and Compute}

% 5 architectures: Mistral-7B-v0.1, OPT-6.7B, GPT-2 XL, Qwen2.5-7B, Pythia-1.4B
% Scaling: Pythia-160M through Pythia-6.9B, Qwen2.5-3B, Phi-3-mini
% Hardware: RunPod A100 80GB
% All inference at fp16, no quantization


%% ============================================================
%% 4. RESULTS
%% ============================================================

\section{Results}

% ────────────────────────────────────────
\subsection{Mode Atlas: $\RV$ Spectral Fingerprint}
% ────────────────────────────────────────

% FIGURE: fig1_mode_atlas_rv — bar chart of R_V across 10 modes
% Key finding: self-referential mode has lowest R_V (strongest contraction)
% d = -1.67 overall vs other modes

% FIGURE: fig6_rv_distribution — violin plot self-ref vs all other modes

% FIGURE: fig5_mode_pairwise_heatmap — pairwise Cohen's d between modes

% FIGURE: fig12_multi_metric_radar — multi-metric discriminant profile


% ────────────────────────────────────────
\subsection{Cross-Architecture Universality}
% ────────────────────────────────────────

% FIGURE: fig2_cross_architecture — forest plot of d across 5 models
% Key finding: replicates in 4/5 architectures at p<0.05
%   Mistral d=-2.26, OPT d=-1.84, GPT-2 XL d=-1.14, Qwen d=-0.72
%   Pythia-1.4B d=-0.31 (NS after FDR)

% TABLE: effect_size_table.tex — comprehensive effect sizes with power


% ────────────────────────────────────────
\subsection{Scaling and Emergence}
% ────────────────────────────────────────

% FIGURE: fig_scaling_curve — effect size vs model parameters
% Key finding: effect present across scales but R²=0.047 (needs more data points)

% FIGURE: fig_training_checkpoints — emergence during training (Pythia-1.4B)
% Key finding: effect emerges early and persists


% ────────────────────────────────────────
\subsection{Causal Evidence: Necessity and Sufficiency}
% ────────────────────────────────────────

% FIGURE: fig8_necessity_sufficiency — 2x2 schematic
% Necessity: dual-layer break d=3.29 (BT+ART drops from 56% to 27.7%)
% Sufficiency: KV patching OR=13.96 (BT+ART uplift d=-3.50)
% Within-session bridge: d=-0.71


% ────────────────────────────────────────
\subsection{Circuit-Level Decomposition}
% ────────────────────────────────────────

% FIGURE: fig_full_head_sweep_32x32 — 1024-head heatmap
% Key finding: 606/1024 heads significant, top L10H20 d=3.90
% Middle layers (L8-L14) show strongest effects — surprising

% SVD per-head: L27_H2 rank deficiency d=-1.54
% Vocabulary projections: top singular directions decode to self-referential tokens


% ────────────────────────────────────────
\subsection{Representation Analysis}
% ────────────────────────────────────────

% Linear probe: 100% accuracy L4+ (n=20, acknowledge overfit risk)
% RSA: 10 modes × 10 layers, self-referential distance trajectory

% FIGURE: fig9_spectral_scatter — PR_early vs PR_late scatter


% ────────────────────────────────────────
\subsection{Statistical Hardening}
% ────────────────────────────────────────

% FIGURE: fig3_statistical_hardening — forest plot with BF
% FIGURE: fig_fdr_correction — FDR summary

% FDR: 30/36 survive BH α=0.05
% Cluster-robust SEs: 10/13 survive (ICC measured for circularity controls)
% Perplexity re-pairing: d=-1.80, p=9.12e-11 (confound ruled out)
% Power: 8/12 effects adequately powered (≥0.80)

% FIGURE: fig10_circularity_controls — structure vs semantics


% ────────────────────────────────────────
\subsection{Safety Applications}
% ────────────────────────────────────────

% FIGURE: fig_safety_roc — ROC curve (AUROC=0.909)
% FIGURE: fig_safety_genuine_vs_deceptive — content vs intent

% Key finding 1: R_V achieves AUROC=0.909 for self-ref detection
% Key finding 2: genuine vs deceptive d=-0.06 (indistinguishable)
% Key finding 3: alignment-faking d=-2.06 vs baseline
% Implication: R_V tracks content, not intent — useful for monitoring but not deception detection


%% ============================================================
%% 5. DISCUSSION
%% ============================================================

\section{Discussion}

% What R_V reveals about transformer geometry
% Why contraction? Information compression hypothesis
% Middle-layer head sweep surprise (vs late-layer R_V effect)
% Content vs intent: implications for AI safety
% The self-feeding loop is negative (d=-4.28): scaffolded > self-feeding
% Homeostatic dynamics: networks resist perturbation

\subsection{Limitations}

% 1. Scaling fit weak (R²=0.047, 6 points)
% 2. Linear probe n=20 (overfit risk)
% 3. Pythia-2.8B checkpoints bugged
% 4. No behavioral validation of "self-reference"
% 5. Template dependence (addressed by cluster-robust SEs)
% 6. Single prompt bank (CANONICAL_CODE/n300)

\subsection{Future Work}

% 1. SAE feature analysis (needs Gemma-2-2B access)
% 2. Multi-seed replication
% 3. Larger scaling study (10+ model sizes)
% 4. Behavioral correlation with self-report accuracy
% 5. Real-time monitoring deployment


%% ============================================================
%% 6. CONCLUSION
%% ============================================================

\section{Conclusion}

% 3-4 sentences summarizing the core contribution

\textbf{TODO: Write conclusion.}


%% ============================================================
%% REFERENCES
%% ============================================================

\bibliographystyle{plainnat}
\bibliography{references}


%% ============================================================
%% APPENDIX
%% ============================================================

\appendix

\section{Complete Effect Size Table}
\label{app:effects}
% \input{effect_size_table.tex}

\section{FDR Correction Details}
\label{app:fdr}
% Full table of 36 tests with original p, q-value, significance

\section{Cluster-Robust Standard Error Details}
\label{app:cluster}
% ICC computation, DEFF, effective n for each experiment

\section{Power Analysis}
\label{app:power}
% Achieved power for each effect, required n for 80% power

\section{Prompt Templates}
\label{app:prompts}
% Example prompts for each of the 10 modes
% Template structure documentation

\section{Full Head Sweep Results}
\label{app:heads}
% Top-50 heads table
% Distribution of d values across all 1024 heads

\end{document}
